% Date : 07/01/2023
% Version : 1
% Auteur : Gaelle Dumas
% Relu : 
% Relecteur : 

% ------------------------------------------------------------ %
% ----------------------  Fiche MCA01  ----------------------- %
% Thème : Cinematique
% Classes : MPSI - PCSI - MP2I
% ------------------------------------------------------------ %



% ======================================================= % 
% ============== Gestion de la compilation ============== %
% ======================================================= % 
\ifdefined\mainIsLoaded\else
\RequirePackage{import}
\subimport{../../}{_preambule_CdE_PC}
\begin{document}
\fi
% ======================================================= % 
% ======================================================= % 






% ============================================================================ %
%                            FICHE D'ENTRAÎNEMENT                              %
% ============================================================================ %
% ======================= Métadonnées sur la fiche =========================== %
\titreFicheEntrainement		{Cinématique}
\grandTheme					{Mécanique}
\numeroFiche				{MCA01}
\uniqueID					{FHGlOcemRQc} % une chaîne de caractères (a-z, A-Z) aléatoire de 10 caractères.
\nombreColonnesReponses		{2} % 1, 2, 3, etc.
% ============================================================================ %



% ===================== Couleurs ======================== %
% ======================================================= % 

% La commande \couleurNBouCouleur prend deux paramètres :
% #1 est la couleur NB ; #2 est la couleur "en couleur"
\def\JRLorange{\couleurNBouCouleur{gray}{orange}}
\def\JRLgreen{\couleurNBouCouleur{gray}{green!70!blue}}
\def\JRLorangeOrBlack{\couleurNBouCouleur{black}{orange}}
\def\JRLgreenOrBlack{\couleurNBouCouleur{black}{green!70!blue}}
\def\JRLblueCC{\couleurNBouCouleur{lightgray!10}{blue!8}}
\def\JRLblue{\couleurNBouCouleur{black}{blue!75}}
\def\JRLblueC{\couleurNBouCouleur{lightgray!30}{blue!15}}
\def\JRLcyan{\couleurNBouCouleur{gray}{cyan}}
\def\JRLcyanC{\couleurNBouCouleur{lightgray!30}{cyan!10}}
\def\JRLred{\couleurNBouCouleur{gray}{red!90!green}}
 \def\JRLblueF{\couleurNBouCouleur{gray}{blue!50}}


%%%%%%%%%%%%%%%%%%%%%%%%%%%%%%%%%%%%%%%%%%%%%%%%%%%%%%%%%%%%%%%%%%%%%%%%%%%%%%%%%%%%%%%%%%%%%%
\debutFicheEntrainement                                                                      %
%%%%%%%%%%%%%%%%%%%%%%%%%%%%%%%%%%%%%%%%%%%%%%%%%%%%%%%%%%%%%%%%%%%%%%%%%%%%%%%%%%%%%%%%%%%%%%




% --------------------------------------------------------------------- %
%                              Prérequis                                %
%                             (facultatif)                              %
% --------------------------------------------------------------------- %
% ================== Métadonnées sur le prérequis ===================== %
\avecPrerequis{Y} % Y ou N
% ===================================================================== %
\begin{prerequis}
	Produit scalaire. Équations différentielles d'ordre 1. Projections de 
	vecteurs.
\end{prerequis}
% --------------------------------------------------------------------- %





%••••••••••••••••••••••••••••••••••••••••••••••••••••••••••••••••••••••••••••
%••••••••••••••••••••••••••••••••••••••••••••••••••••••••••••••••••••••••••••
\sectionFicheEntrainement{Déplacements rectilignes}
%••••••••••••••••••••••••••••••••••••••••••••••••••••••••••••••••••••••••••••
%••••••••••••••••••••••••••••••••••••••••••••••••••••••••••••••••••••••••••••







% ********************************************************************* %
%                           ENTRAÎNEMENT                                % 
% ********************************************************************* %
% ================ Métadonnées sur l'entraînement ===================== %

\titreEntrainementFacultatif	{Distance et temps de parcours}
\hauteurLargeurCadreReponse		{7mm}{3.5cm}
\dureeResolutionFacultative		{1} % 1, 2, 3 ou 4
\typeEntrainement				{Newton} % Newton ou QCM ou AN ou vide (exercice "normal")
\basiqueEtTransversal			{Y}
\nombreColonnesQuestions		{1} % vide, 1, 2, 3, etc.
\avecPlusieursQuestions			{Y} % Y ou N
\initialisationEntrainement
% ===================================================================== %

Une voiture se déplace en ligne droite à $\SI{90}{km.h^{-1}}$.

{\itshape Toutes les réponses seront exprimées en \glm{heures-minutes-secondes}, par exemple \glm{$\SI{2}{h}\ \SI{32}{min}\ \SI{12}{s}$}.} 

% =============================================================================== %
\debutEntrainement
% =============================================================================== %

% ^^^^^^^^^^^^^^^^^^^^^^^^^^^^^^^^^^^^^^ %
%               Question                 %
% ^^^^^^^^^^^^^^^^^^^^^^^^^^^^^^^^^^^^^^ %

\begin{enonce}
Combien de temps faut-il à cette voiture pour parcourir $\SI{100}{km}$ ?
\end{enonce}

\reponse{$\SI{1}{h}\ \SI{6}{min}\ \SI{40}{s}$}

\begin{corrige}
La voiture avance à vitesse constante.
Pour parcourir $\SI{100}{km}$, il lui faudra le temps 
\[
\tau=\frac{\SI{100}{km}}{\SI{90}{km.h^{-1}}}=\SI{1,11}{h}=\SI{1}{h} \ \SI{6}{min}\ \SI{40}{s}.
\] 
\end{corrige}

% ************************************** %



% ^^^^^^^^^^^^^^^^^^^^^^^^^^^^^^^^^^^^^^ %
%               Question                 %
% ^^^^^^^^^^^^^^^^^^^^^^^^^^^^^^^^^^^^^^ %

\begin{enonce}
Quel serait l'allongement du temps de trajet si elle roulait à $\SI{80}{km.h^{-1}}$ ?
\end{enonce}

\reponse{$\SI{8}{min}\ \SI{20}{s}$}

\begin{corrige}
	Pour parcourir $\SI{100}{km}$ à $\SI{80}{km.h^{-1}}$, il lui faudrait le temps $\tau'=\frac{\SI{100}{km}}{\SI{80}{km.h^{-1}}}=\SI{1,25}{h}$.
	Le temps de trajet serait donc allongé de $\Delta t=\tau'-\tau=\SI{0,14}{h}= \SI{8}{min}\ \SI{20}{s}$. 
\end{corrige}

% ************************************** %

% =============================================================================== %
\finEntrainement
% =============================================================================== %







% ********************************************************************* %
%                           ENTRAÎNEMENT                                % 
% ********************************************************************* %
% ================ Métadonnées sur l'entraînement ===================== %

\titreEntrainementFacultatif	{Distance parcourue}
\hauteurLargeurCadreReponse		{7mm}{3.5cm}
\dureeResolutionFacultative		{2} % 1, 2, 3 ou 4
\typeEntrainement				{Newton} % Newton ou QCM ou AN ou vide (exercice "normal")
\basiqueEtTransversal			{Y}
\nombreColonnesQuestions		{1} % vide, 1, 2, 3, etc.
\avecPlusieursQuestions			{Y} % Y ou N
\initialisationEntrainement
% ===================================================================== %

Une voiture se déplace en ligne droite. 
Initialement à l'arrêt, elle subit une accélération constante valant $a_0$
 pendant une durée $\tau_1$, puis continue à vitesse constante pendant une durée $\tau_2$.

% =============================================================================== %
\debutEntrainement
% =============================================================================== %

% ^^^^^^^^^^^^^^^^^^^^^^^^^^^^^^^^^^^^^^ %
%               Question                 %
% ^^^^^^^^^^^^^^^^^^^^^^^^^^^^^^^^^^^^^^ %

\begin{enonce}
	Quelle est la vitesse $v_1$ du véhicule à la date $t=\tau_1$ ?
\end{enonce}

\reponse{$a_0 \times \tau_1$}

\begin{corrige}
L'accélération est constante durant le temps $\tau_1$ et la vitesse initiale est nulle. 
La vitesse à un instant $t$ vaut donc $v(t) = a_0 \times t$, d'où 
$v_1 = v(\tau_1) = a_0 \times \tau_1$.
\end{corrige}

% ************************************** %


% ^^^^^^^^^^^^^^^^^^^^^^^^^^^^^^^^^^^^^^ %
%               Question                 %
% ^^^^^^^^^^^^^^^^^^^^^^^^^^^^^^^^^^^^^^ %

\begin{enonce}
	Quelle est la distance parcourue durant $\tau_1$ ?
\end{enonce}

\reponse{$\dfrac{a_0\times {\tau_1}^2}{2}$}

\begin{corrige}
Pour $t\in[0, \tau_1]$, la vitesse est décrite par l'équation : $v(t) = a_0\times t$. 
La distance parcourue à la date $t$, s'écrit donc $d(t)=\frac{1}{2}a_0\times t^2$. Ainsi, on a $d_1=d(\tau_1)=\frac{a_0\times{\tau_1}^2}{2}$.
\end{corrige}

% ************************************** %


% ^^^^^^^^^^^^^^^^^^^^^^^^^^^^^^^^^^^^^^ %
%               Question                 %
% ^^^^^^^^^^^^^^^^^^^^^^^^^^^^^^^^^^^^^^ %

\begin{enonce}
	Quelle est la distance totale parcourue en fonction de $a_0$, $\tau_1$ et $\tau_2$ ?
\end{enonce}

\reponse{$a_0\times \tau_1 \times \left(\dfrac{\tau_1}{2} + \tau_2\right)$}

\begin{corrige}
	La distante totale parcourue est $d_{\text{tot}} = d_1+d_2$ avec  $d_1$ évaluée à la question précédente et $d_2$ la distance parcourue par le véhicule dans la seconde phase du mouvement où il progresse à vitesse constante. 
	
	Or, on a  $d_2 = v_1 \times \tau_2$. Ainsi, on a $d_{\text{tot}} = 
	a_0 \times \tau_1 \times\left(\dfrac{\tau_1}{2} + \tau_2\right)$.
\end{corrige}

% ************************************** %

% =============================================================================== %
\finEntrainement
% =============================================================================== %



% ********************************************************************* %
%                           ENTRAÎNEMENT                                % 
% ********************************************************************* %
% ================ Métadonnées sur l'entraînement ===================== %

\titreEntrainementFacultatif	{Longueur d'une piste de décollage}
\hauteurLargeurCadreReponse		{7mm}{1.5cm}
\dureeResolutionFacultative		{3} % 1, 2, 3 ou 4
\typeEntrainement				{QCM} % Newton ou QCM ou AN ou vide (exercice "normal")
\basiqueEtTransversal			{Y}
\nombreColonnesQuestions		{1} % vide, 1, 2, 3, etc.
\avecPlusieursQuestions			{N} % Y ou N
\initialisationEntrainement
% ===================================================================== %


% =============================================================================== %
\debutEntrainement
% =============================================================================== %

% ^^^^^^^^^^^^^^^^^^^^^^^^^^^^^^^^^^^^^^ %
%               Question                 %
% ^^^^^^^^^^^^^^^^^^^^^^^^^^^^^^^^^^^^^^ %

\begin{enonce}
Pour décoller un avion doit atteindre la vitesse de $v_d= \SI{180}{km.h^{-1}}$ en bout de piste.\newline
Quelle est la longueur minimale $L$ de la piste de décollage si l'avion 
accélère  uniformément à la valeur $a=\SI{2,5}{m.s^{-2}}$ ?
\begin{listeQCM4Colonnes}
	\item \np[m]{300} 
	\item \np[m]{450} 
	\item \np[m]{500}  
	\item \np[m]{650}  
\end{listeQCM4Colonnes}
\medskip
\end{enonce}

\reponse{\reponseC{}}

\begin{corrige}
À $t=0$, l'avion a une vitesse nulle. Sa vitesse au temps $t$ s'écrit alors  
$v(t) = a \times t $ et la distance qu'il parcourt, vaut $d(t) = \dfrac{1}{2}a\times t^2$.

D'abord le temps $t_d$ où l'avion atteint la vitesse $v_d$ vaut $t_d = \dfrac{v_d}{a}$.

Pour faire l'application numérique, il nous faut  exprimer la vitesse $v_d$ en $\si{m.s^{-1}}$ : on a 
\[
v_d=\frac{\SI{180e3}{m}}{\SI{3600}{s}}=\SI{50}{m.s^{-1}}
\text{ et donc }
t_d=\frac{\SI{50}{m.s^{-1}}}{\SI{2,5}{m.s^{-2}}} = \SI{20}{s}.
\]
La longueur de la piste correspond à la distance parcourue pendant cette durée : c'est
\[
L =\dfrac{1}{2}a\times {t_d}^2 = \dfrac{{v_d}^2}{2a} = \dfrac{(\SI{50}{m.s^{-1}})^2}{2 \times \SI{2,5}{m.s^{-2}}} =\SI{500}{m}.
\]
\end{corrige}

% ************************************** %

% =============================================================================== %
\finEntrainement
% =============================================================================== %





% ********************************************************************* %
%                           ENTRAÎNEMENT                                % 
% ********************************************************************* %
% ================ Métadonnées sur l'entraînement ===================== %

\titreEntrainementFacultatif	{Distance de freinage}
\hauteurLargeurCadreReponse		{7mm}{1.5cm}
\dureeResolutionFacultative		{1} % 1, 2, 3 ou 4
\typeEntrainement				{QCM} % Newton ou QCM ou AN ou vide (exercice "normal")
\basiqueEtTransversal			{Y}
\nombreColonnesQuestions		{1} % vide, 1, 2, 3, etc.
\avecPlusieursQuestions			{N} % Y ou N
\initialisationEntrainement
% ===================================================================== %


% =============================================================================== %
\debutEntrainement
% =============================================================================== %

% ^^^^^^^^^^^^^^^^^^^^^^^^^^^^^^^^^^^^^^ %
%               Question                 %
% ^^^^^^^^^^^^^^^^^^^^^^^^^^^^^^^^^^^^^^ %

\begin{enonce}
Une voiture roule à $\SI{110}{km.h^{-1}}$ en ligne droite. En supposant que les 
freins imposent une décélération constante de norme $a=\SI{10}{m.s^{-2}}$, 
déterminer la distance d'arrêt de la voiture. 
	\begin{listeQCM4Colonnes}
		\item $\SI{37,8}{m}$ 
		\item $\SI{46,7}{m}$ 
		\item $\SI{55,9}{m}$
		\item $\SI{63,5}{m}$  
	\end{listeQCM4Colonnes}
	\medskip
\end{enonce}

\reponse{\reponseB{}}

\begin{corrigeNewpage}
	La vitesse de la voiture à un instant $t$ s'écrit $v(t) = v_i - a \times t$ avec 
	\[
	v_i = \SI{110}{km.h^{-1}}=\frac{\SI{110e3}{m}}{\SI{3600}{s}}= \SI{30,6}{m.s^{-1}}.
	\]
	Ainsi, le véhicule s'arrêtera à la date $t_a$ telle que $v_i - a \times t = \SI{0}{m.s^{-1}}$. On a  $t_a=\dfrac{v_i}{a} =\frac{\SI{30,6}{m.s^{-1}}}{\SI{10}{m.s^{-2}}}=\SI{3,06}{s}$.
	
	La distance parcourue pendant le freinage vaut $d(t) = v_i\times t -\dfrac{1}{2}a \times t^2$. 
	
	La distance d'arrêt 
	$d_a$ correspond à la distance parcourue pendant la durée $t_a$ : c'est $d_a =\frac{{v_i}^2}{2a} =\SI{46,7}{m}$.
\end{corrigeNewpage}

% ************************************** %

% =============================================================================== %
\finEntrainement
% =============================================================================== %



% •••••••••••••••••••••••••••••••••••••••••••••••••••••••••••••••••••••••••••• %
% •••••••••••••••••••••••••••••••••••••••••••••••••••••••••••••••••••••••••••• %
\sectionFicheEntrainement{Coordonnées et projections de vecteurs}
% •••••••••••••••••••••••••••••••••••••••••••••••••••••••••••••••••••••••••••• %
% •••••••••••••••••••••••••••••••••••••••••••••••••••••••••••••••••••••••••••• %






% ********************************************************************* %
%                           ENTRAÎNEMENT                                % 
% ********************************************************************* %
% ================ Métadonnées sur l'entraînement ===================== %

\titreEntrainementFacultatif	{Composantes de vecteurs}
\hauteurLargeurCadreReponse		{8mm}{5cm}
\dureeResolutionFacultative		{2} % 1, 2, 3 ou 4
\typeEntrainement				{AN} % Newton ou QCM ou AN ou vide (exercice "normal")
\basiqueEtTransversal			{Y}
\nombreColonnesQuestions		{1} % vide, 1, 2, 3, etc.
\avecPlusieursQuestions			{Y} % Y ou N
\initialisationEntrainement
% ===================================================================== %


% mmmmmmmmmmmmmmmmmmmmmmmmmmmmmmmmmmmmmmmmmmmmmmmmmmmmmmmmm %
% 		  Insertion d'une image en regard d'un texte
% mmmmmmmmmmmmmmmmmmmmmmmmmmmmmmmmmmmmmmmmmmmmmmmmmmmmmmmmm %
% --------------------------------------------------------- %
\pourcentageDeLaPartieAGauche   {0.7}
% --------------------------------------------------------- %
%                                                           %
% -------------------- Partie à gauche -------------------- %
%                                                           %
\initialisationPartieGauche % 
%                                                           %
%                                                           %
 On considère deux points $\ptA$ et $\ptB$ tels que la droite  $(\ptA\ptB)$ est parallèle à la droite $(\ptO y)$. Le vecteur $\vv{\ptO \ptA}$  fait un 
 angle $\theta$ avec l'axe $(\ptO x)$.
 
 Exprimer les composantes des vecteurs suivants dans le repère  $(\ptO,\vv{e_x},\vv{e_y})$ en fonction de 
 $a=\big\|\vv{\ptO \ptA}\big\|$, $b=\big\|\vv{\ptA \ptB}\big\|$ et de l'angle $\theta$.

% --------------------------------------------------------- %
%                                                           %
% -------------------- Partie à droite -------------------- %
%                                                           %
\initialisationPartieDroite %
%                                                           %
%                                                           %
\begin{center}
	\subimport{_images/}{fiche_MCA1_ABC_cart_GDG_v2.tex}
\end{center}
% --------------------------------------------------------- %
\finalisationDuPartageDePage %					
% --------------------------------------------------------- %


% =============================================================================== %
\debutEntrainement
% =============================================================================== %


% ^^^^^^^^^^^^^^^^^^^^^^^^^^^^^^^^^^^^^^ %
%               Question                 %
% ^^^^^^^^^^^^^^^^^^^^^^^^^^^^^^^^^^^^^^ %

\begin{enonce}
	$\vv{\ptO \ptA}$
\end{enonce}

\reponse{$a \left(\cos(\theta) \vv{e_x}+\sin(\theta) \vv{e_y}\right)$}

\begin{corrige}
	On a $\vv{\ptO \ptA} = a \left(\cos(\theta) \vv{e_x}+\sin(\theta) \vv{e_y}\right)$.
\end{corrige}

% ************************************** %




% ^^^^^^^^^^^^^^^^^^^^^^^^^^^^^^^^^^^^^^ %
%               Question                 %
% ^^^^^^^^^^^^^^^^^^^^^^^^^^^^^^^^^^^^^^ %
\begin{enonce}
	$\vv{\ptO \ptB}$ 
\end{enonce}

\reponse{$a \left(\cos(\theta) \vv{e_x}+\left(\sin(\theta)+\frac{b}{a}\right)\vv{e_y}\right)$}

\begin{corrige}
	On a $\vv{\ptO \ptB} =\vv{\ptO \ptA}+\vv{\ptA \ptB}= a \left(\cos(\theta) \vv{e_x}+\left(\sin(\theta)+\frac{b}{a}\right)\vv{e_y}\right)$.
\end{corrige}

% ************************************** %



% ^^^^^^^^^^^^^^^^^^^^^^^^^^^^^^^^^^^^^^ %
%               Question                 %
% ^^^^^^^^^^^^^^^^^^^^^^^^^^^^^^^^^^^^^^ %

\begin{enonce}
 $\vv{\ptO \ptA} + \vv{\ptO \ptB}$
\end{enonce}

\reponse{$ a \left(2\cos(\theta) \vv{e_x}+\left(2\sin(\theta)+\frac{b}{a}\right)\vv{e_y}\right)$}

\begin{corrige}
	On a $\vv{\ptO \ptA}+ \vv{\ptO \ptB} = a \left(2\cos(\theta) \vv{e_x}+\left(2\sin(\theta)+\frac{b}{a}\right)\vv{e_y}\right)$.
\end{corrige}

% ************************************** %



% ^^^^^^^^^^^^^^^^^^^^^^^^^^^^^^^^^^^^^^ %
%               Question                 %
% ^^^^^^^^^^^^^^^^^^^^^^^^^^^^^^^^^^^^^^ %

\begin{enonce}
  $\vv{\ptO \ptA}-\vv{\ptO \ptB}$
\end{enonce}

\reponse{$-b \vv{e_y}$}

\begin{corrige}
	On a $\vv{\ptO \ptA}-\vv{\ptO \ptB} =\vv{\ptB \ptA}=-b \vv{e_y}$.
\end{corrige}

% ************************************** %



% =============================================================================== %
\finEntrainement
% =============================================================================== %







% ********************************************************************* %
%                           ENTRAÎNEMENT                                % 
% ********************************************************************* %
% ================ Métadonnées sur l'entraînement ===================== %
\titreEntrainementFacultatif	{Les coordonnées cylindriques}
\hauteurLargeurCadreReponse		{8mm}{4cm}
\dureeResolutionFacultative		{2} % 1, 2, 3 ou 4
\typeEntrainement				{Newton} % Newton ou QCM ou AN ou vide (exercice "normal")
\basiqueEtTransversal			{Y}
\nombreColonnesQuestions		{1} % vide, 1, 2, 3, etc.
\avecPlusieursQuestions			{Y} % Y ou N
\initialisationEntrainement
% ===================================================================== %


% mmmmmmmmmmmmmmmmmmmmmmmmmmmmmmmmmmmmmmmmmmmmmmmmmmmmmmmmm %
% 		  Insertion d'une image en regard d'un texte
% mmmmmmmmmmmmmmmmmmmmmmmmmmmmmmmmmmmmmmmmmmmmmmmmmmmmmmmmm %
% --------------------------------------------------------- %
\pourcentageDeLaPartieAGauche   {0.48}
% --------------------------------------------------------- %
%                                                           %
% -------------------- Partie à gauche -------------------- %
%                                                           %
\initialisationPartieGauche % 
%                                                           %
%                                                           %
 On considère le schéma ci-contre, dans lequel 
 \begin{itemize}
 \item la base cartésienne $(\vv{e_x},\vv{e_y},\vv{e_z})$ 
 \item et la base cylindrique  $(\vv{e_r},\vv{e_\theta},\vv{e_z})$
 \end{itemize}
  sont définies.
  
  Le point $\ptM$ est repéré par la donnée de $r$, $\theta$ et $z$.
% --------------------------------------------------------- %
%                                                           %
% -------------------- Partie à droite -------------------- %
%                                                           %
\initialisationPartieDroite %
%                                                           %
%                                                           %
\begin{center}
	\subimport{_images/}{fiche_sup_coordonnees_cylindrique_CBD_v2.tex}
\end{center}
% --------------------------------------------------------- %
\finalisationDuPartageDePage %					
% --------------------------------------------------------- %




% =============================================================================== %
\debutEntrainement
% =============================================================================== %


% ^^^^^^^^^^^^^^^^^^^^^^^^^^^^^^^^^^^^^^ %
%               Question                 %
% ^^^^^^^^^^^^^^^^^^^^^^^^^^^^^^^^^^^^^^ %

\begin{enonce}
	Écrire le vecteur $\vv{\ptO \ptM'}$ dans la base cartésienne
\end{enonce}

\reponse{$r \left(\cos(\theta) \vv{e_x}+\sin(\theta) \vv{e_y}\right)$}

\begin{corrige}
	On a $\vv{\ptO \ptM'}=r \left(\cos(\theta) \vv{e_x}+\sin(\theta) \vv{e_y}\right)$.
\end{corrige}

% ************************************** %


% ^^^^^^^^^^^^^^^^^^^^^^^^^^^^^^^^^^^^^^ %
%               Question                 %
% ^^^^^^^^^^^^^^^^^^^^^^^^^^^^^^^^^^^^^^ %

\begin{enonce}
	Écrire le vecteur $\vv{\ptO \ptM'}$ dans la base cylindrique
\end{enonce}

\reponse{$r \vv{e_r}$}

\begin{corrige}
	On a $\vv{\ptO \ptM'}=r \vv{e_r}$.
\end{corrige}

% ************************************** %


% ^^^^^^^^^^^^^^^^^^^^^^^^^^^^^^^^^^^^^^ %
%               Question                 %
% ^^^^^^^^^^^^^^^^^^^^^^^^^^^^^^^^^^^^^^ %

\begin{enonce}
Écrire le vecteur $\vv{\ptO \ptM}$ dans la base cartésienne
\end{enonce}

\reponse{$r \left(\cos(\theta) \vv{e_x}+\sin(\theta) \vv{e_y}\right) +z\vv{e_z} $}

\begin{corrige}
	On a $\vv{\ptO \ptM}=r \left(\cos(\theta) \vv{e_x}+\sin(\theta) \vv{e_y}\right) +z\vv{e_z}$.
\end{corrige}

% ************************************** %


% ^^^^^^^^^^^^^^^^^^^^^^^^^^^^^^^^^^^^^^ %
%               Question                 %
% ^^^^^^^^^^^^^^^^^^^^^^^^^^^^^^^^^^^^^^ %

\begin{enonce}
	Écrire le vecteur $\vv{\ptO \ptM}$ dans la base cylindrique 
\end{enonce}

\reponse{$r \vv{e_r}+z\vv{e_z}$}

\begin{corrige}
	On a $\vv{\ptO \ptM}=r \vv{e_r}+z\vv{e_z}$.
\end{corrige}

% ************************************** %

% =============================================================================== %
\finEntrainement
% =============================================================================== %




\newpage


% ********************************************************************* %
%                           ENTRAÎNEMENT                                % 
% ********************************************************************* %
% ================ Métadonnées sur l'entraînement ===================== %
\titreEntrainementFacultatif	{Les coordonnées sphériques}
\hauteurLargeurCadreReponse		{8mm}{4cm}
\dureeResolutionFacultative		{2} % 1, 2, 3 ou 4
\typeEntrainement				{Newton} % Newton ou QCM ou AN ou vide (exercice "normal")
\basiqueEtTransversal			{Y}
\nombreColonnesQuestions		{1} % vide, 1, 2, 3, etc.
\avecPlusieursQuestions			{Y} % Y ou N
\initialisationEntrainement
% ===================================================================== %


 On considère le schéma ci-dessous, dans lequel la base cartésienne $(\vv{e_x},\vv{e_y},\vv{e_z})$ 
et la base sphérique $(\vv{e_r},\vv{e_\theta},\vv{e_\varphi})$
  sont définies.


\begin{center}
	\subimport{_images/}{fiche_sup_coordonnees_spheriques_ECA.tex}
\end{center}


  Le point $\ptM$ est repéré par la donnée de $r$, $\theta$ et $\varphi$.


% =============================================================================== %
\debutEntrainement
% =============================================================================== %


% ^^^^^^^^^^^^^^^^^^^^^^^^^^^^^^^^^^^^^^ %
%               Question                 %
% ^^^^^^^^^^^^^^^^^^^^^^^^^^^^^^^^^^^^^^ %

\begin{enonce}
	Écrire la norme de $\vv{\ptO \ptM'}$ en fonction de $r$ et $\theta$
\end{enonce}

\reponse{$\left|r\sin(\theta)\right|$}

\begin{corrige}
	On a $\big\|\vv{\ptO \ptM'}\big\|=\left|r\sin(\theta)\right|$.
\end{corrige}

% ************************************** %


% ^^^^^^^^^^^^^^^^^^^^^^^^^^^^^^^^^^^^^^ %
%               Question                 %
% ^^^^^^^^^^^^^^^^^^^^^^^^^^^^^^^^^^^^^^ %
\begin{enonce}
	Écrire le vecteur $\vv{\ptO \ptM'}$ dans la base cartésienne
\end{enonce}

\reponse{$r\sin(\theta)  \left(\cos(\varphi) \vv{e_x}+\sin(\varphi) \vv{e_y}\right)$}

\begin{corrige}
	On a $\vv{\ptO \ptM'}=r\sin(\theta)  \left(\cos(\varphi) \vv{e_x}+\sin(\varphi) \vv{e_y}\right)$.
\end{corrige}

% ************************************** %


% ^^^^^^^^^^^^^^^^^^^^^^^^^^^^^^^^^^^^^^ %
%               Question                 %
% ^^^^^^^^^^^^^^^^^^^^^^^^^^^^^^^^^^^^^^ %

\begin{enonce}
	Écrire le vecteur $\vv{\ptO \ptM}$ dans la base cartésienne
	\end{enonce}
	
	\reponse{\small$r\sin(\theta)  \left(\cos(\varphi) \vv{e_x}+\sin(\varphi)\vv{e_y}\right)+r\cos(\theta)\vv{e_z}$}
	
	\begin{corrige}
		On a $\vv{\ptO \ptM}=\vv{\ptO \ptM'}+\vv{\ptM' \ptM}=r\sin(\theta)  \left(\cos(\varphi) \vv{e_x}+\sin(\varphi) \vv{e_y}\right)+r\cos(\theta)\vv{e_z}$.
	\end{corrige}
	
	% ************************************** %
	
	
	% ^^^^^^^^^^^^^^^^^^^^^^^^^^^^^^^^^^^^^^ %
	%               Question                 %
	% ^^^^^^^^^^^^^^^^^^^^^^^^^^^^^^^^^^^^^^ %
	
	\begin{enonce}
		Écrire le vecteur $\vv{\ptO \ptM}$ dans la base sphérique
	\end{enonce}
	
	\reponse{$r \vv{e_r} $}
	
	\begin{corrige}
		On a $\vv{\ptO \ptM}=r \vv{e_r}$.
	\end{corrige}
	
	% ************************************** %
	
	
	% ^^^^^^^^^^^^^^^^^^^^^^^^^^^^^^^^^^^^^^ %
	%               Question                 %
	% ^^^^^^^^^^^^^^^^^^^^^^^^^^^^^^^^^^^^^^ %
	
	\begin{enonce}
		Écrire le vecteur $\vv{e_z}$ dans la base sphérique
	\end{enonce}

\reponse{$\cos(\theta)\,\vv{e_r}-\sin(\theta)\,\vv{e_\theta}$}

\begin{corrige}
Calculons les projections de $\vv{e_z}$ sur les trois vecteurs de la base sphérique : on a
\[
\begin{cases}
\vv{e_z}\cdot\vv{e_r}=\cos(\theta) \\
\vv{e_z}\cdot\vv{e_\theta}=\cos\left(\theta+\frac{\pi}{2}\right)=-\sin(\theta)\\% \qquad\text{et}\qquad
\vv{e_z}\cdot\vv{e_\varphi}=0.
\end{cases}
\]
Par conséquent, on a
\[
\vv{e_z}=\left(\vv{e_z}\cdot\vv{e_r}\right)\,\vv{e_r}+
\left(\vv{e_z}\cdot\vv{e_\theta}\right)\,\vv{e_\theta}+
\left(\vv{e_z}\cdot\vv{e_\varphi}\right)\,\vv{e_\varphi}=
\cos(\theta)\,\vv{e_r}-\sin(\theta)\,\vv{e_\theta}.
\]
\end{corrige}

% ************************************** %

% =============================================================================== %
\finEntrainement
% =============================================================================== %






% ********************************************************************* %
%                           ENTRAÎNEMENT                                % 
% ********************************************************************* %
% ================ Métadonnées sur l'entraînement ===================== %
\titreEntrainementFacultatif	{Jouons au tennis}
\hauteurLargeurCadreReponse		{8mm}{4cm}
\dureeResolutionFacultative		{2} % 1, 2, 3 ou 4
\typeEntrainement				{Newton} % Newton ou QCM ou AN ou vide (exercice "normal")
\basiqueEtTransversal			{Y}
\nombreColonnesQuestions		{1} % vide, 1, 2, 3, etc.
\avecPlusieursQuestions			{Y} % Y ou N
\initialisationEntrainement
% ===================================================================== %

Un élève regarde un match de tennis. Il filme un des échanges et décide d'étudier le mouvement de la balle pour en déduire sa vitesse et son accélération.
 
Pour cela, il utilise un logiciel d'exploitation de vidéo et remplit le tableau suivant :
\begin{center}
\renewcommand*{\arraystretch}{1.2}
	\begin{tabular}{|c|c|c|c|c|c|}
		\hline
		$t$ (en \si{s})	& $0$ & $\np{0,05}$ & $\np{0,10}$ & $\np{0,15}$  & $\np{0,20}$ \\
		\hline
		$x$ (en \si{m}) & $0$ & $\np{0,35}$ & $\np{0,70}$ & $\np{1,05}$ & $\np{1,40}$ \\
		\hline
		$y$ (en \si{m}) & $\np{1,5}$ & $\np{2,09}$ & $\np{2,66}$ & $\np{3,21}$ & $\np{3,74}$ \\
		\hline
	\end{tabular}
\renewcommand*{\arraystretch}{1}
\end{center}

% =============================================================================== %
\debutEntrainement
% =============================================================================== %


% ^^^^^^^^^^^^^^^^^^^^^^^^^^^^^^^^^^^^^^ %
%               Question                 %
% ^^^^^^^^^^^^^^^^^^^^^^^^^^^^^^^^^^^^^^ %
\begin{enonce}
	Déterminer la vitesse $v_0$ (en $\si{km.h^{-1}}$) de la balle à l'instant initial 
\end{enonce}

\reponse{\SI{49,4}{km.h^{-1}}}

\begin{corrigeNewpage}
	La vitesse de la balle %repérée par le point $\ptM$ 
	à l'instant $t_1$, s'écrit $\vv{v}(\ptM,t_1)=v_x(t_1) \vv{e_x} + v_y(t_1) \vv{e_y}$ 
	avec  
	\[
	v_x(t_1)\simeq\frac{x(t_1+\Delta t)-x(t_1)}{ \Delta t}, \quad 
	v_y(t_1)\simeq\frac{y(t_1+\Delta t)-y(t_1)}{ \Delta t}  \quad\text{et}\quad \Delta t=\SI{0,05}{s}.
	\] 
	Nous obtenons le tableau suivant : \minusSmallskip
\begin{center}
\renewcommand*{\arraystretch}{1.2}
	\begin{tabular}{|c|c|c|c|c|}
		\hline
		$t$ (en \si{s})	& $0$ & $\np{0,05}$ & $\np{0,10}$ & $\np{0,15}$   \\
		\hline
		$v_x$ (en \si{m.s^{-1}})  & $7$ & $7$ & $7$& $7$\\
		\hline
		$v_y$ (en \si{m.s^{-1}})  & $\np{11,8}$ & $\np{11,4}$ & $\np{11.0}$ & $\np{10,6}$ \\
		\hline
	\end{tabular}
\renewcommand*{\arraystretch}{1}
\end{center}
À l'instant initial, nous pouvons écrire : $v_0 \simeq \sqrt{\big(\SI{7}{m.s^{-1}}\big)^2+\big(\SI{11,8}{m.s^{-1}}\big)^2}=\SI{13,72}{m.s^{-1}}=\SI{49,4}{km.h^{-1}}$.
\end{corrigeNewpage}

% ************************************** %


% ^^^^^^^^^^^^^^^^^^^^^^^^^^^^^^^^^^^^^^ %
%               Question                 %
% ^^^^^^^^^^^^^^^^^^^^^^^^^^^^^^^^^^^^^^ %

\begin{enonce}
	Déterminer l'accélération (en $\si{m.s^{-2}}$) de la balle à l'instant initial 
\end{enonce}

\reponse{\SI{8.0}{m.s^{-2}}}

\begin{corrige}
	L'accélération de la balle à l'instant $t_1$, s'écrit $\vv{a}(\ptM,t_1)=a_x(t_1) \vv{e_x} + a_y(t_1) \vv{e_y}$
	avec  
	\[
	a_x(t_1)\simeq\frac{v_x(t_1+\Delta t)-v_x(t_1)}{ \Delta t}, \quad a_y(t_1)\simeq\frac{v_y(t_1+\Delta t)-v_y(t_1)}{ \Delta t}
	\quad\text{et}\quad \Delta t=\SI{0,05}{s}.
	\]
	ce qui donne 
	\[
	a_x(0)\simeq
	\frac{\SI{7}{m.s^{-1}}-\SI{7}{m.s^{-1}}}{\SI{.05}{s}}=
	\SI{0}{m.s^{-2}}
	\quad\text{et}\quad
	a_y(0)\simeq
	\frac{\SI{11.4}{m.s^{-1}}-\SI{11.8}{m.s^{-1}}}{\SI{.05}{s}}=
	\SI{-8}{m.s^{-2}}
	\]
	L'accélération initiale vaut donc $a_0 \simeq \sqrt{\big(\SI{0}{m.s^{-2}}\big)^2+\big(\SI{-8}{m.s^{-2}}\big)^2}=\SI{8.0}{m.s^{-2}}$.
\end{corrige}

% ************************************** %

% =============================================================================== %
\finEntrainement
% =============================================================================== %




\newpage



% •••••••••••••••••••••••••••••••••••••••••••••••••••••••••••••••••••••••••••• %
% •••••••••••••••••••••••••••••••••••••••••••••••••••••••••••••••••••••••••••• %
\sectionFicheEntrainement{Dérivée de vecteurs}
% •••••••••••••••••••••••••••••••••••••••••••••••••••••••••••••••••••••••••••• %
% •••••••••••••••••••••••••••••••••••••••••••••••••••••••••••••••••••••••••••• %






% ********************************************************************* %
%                           ENTRAÎNEMENT                                % 
% ********************************************************************* %
% ================ Métadonnées sur l'entraînement ===================== %

\titreEntrainementFacultatif	{Étude d'un mouvement hélicoïdal}
\hauteurLargeurCadreReponse		{8mm}{5cm}
\dureeResolutionFacultative		{2} % 1, 2, 3 ou 4
\typeEntrainement				{} % Newton ou QCM ou AN ou vide (exercice "normal")
\basiqueEtTransversal			{Y}
\nombreColonnesQuestions		{1} % vide, 1, 2, 3, etc.
\avecPlusieursQuestions			{Y} % Y ou N
\initialisationEntrainement
% ===================================================================== %


% mmmmmmmmmmmmmmmmmmmmmmmmmmmmmmmmmmmmmmmmmmmmmmmmmmmmmmmmm %
% 		  Insertion d'une image en regard d'un texte
% mmmmmmmmmmmmmmmmmmmmmmmmmmmmmmmmmmmmmmmmmmmmmmmmmmmmmmmmm %
% --------------------------------------------------------- %
\pourcentageDeLaPartieAGauche   {0.48}
% --------------------------------------------------------- %
%                                                           %
% -------------------- Partie à gauche -------------------- %
%                                                           %
\initialisationPartieGauche % 
%                                                           %
%                                                           %
Le point matériel $\ptM$ de coordonnées cartésiennes $(x, y, z)$ décrit une trajectoire hélicoïdale, définie par les équations : 
\[
\begin{cases}
x(t)=a\times \cos(\omega t) \\
y(t)=a\times \sin(\omega t) \\
z(t)=b\times t.
\end{cases}
\] 
% --------------------------------------------------------- %
%                                                           %
% -------------------- Partie à droite -------------------- %
%                                                           %
\initialisationPartieDroite %
%                                                           %
%                                                           %
\begin{center}
	\subimport{_images/}{fiche_sup_1.10_ECA.tex}
\end{center}
% --------------------------------------------------------- %
\finalisationDuPartageDePage %					
% --------------------------------------------------------- %




% =============================================================================== %
\debutEntrainement
% =============================================================================== %


% ^^^^^^^^^^^^^^^^^^^^^^^^^^^^^^^^^^^^^^ %
%               Question                 %
% ^^^^^^^^^^^^^^^^^^^^^^^^^^^^^^^^^^^^^^ %

\begin{enonce}
	Déterminer la vitesse $\vv{v}(\ptM)$ dans la base cartésienne 
\end{enonce}

\reponse{$a\omega\left(-\sin(\omega t) \vv{e_x}+\cos(\omega t) \vv{e_y}\right)+b\vv{e_z}$}

\begin{corrige}
	On a $\vv{v}(\ptM)=\dot{x}\, \vv{e_x}+\dot{y}\, \vv{e_y}+\dot{z}\, \vv{e_z}=a\omega\big(-\sin(\omega t) \vv{e_x}+\cos(\omega t) \vv{e_y}\big)+b\vv{e_z}$.
\end{corrige}

% ************************************** %


% ^^^^^^^^^^^^^^^^^^^^^^^^^^^^^^^^^^^^^^ %
%               Question                 %
% ^^^^^^^^^^^^^^^^^^^^^^^^^^^^^^^^^^^^^^ %

\begin{enonce}
	Déterminer la norme de la vitesse
\end{enonce}

\reponse{$\sqrt{\left(a\omega\right)^2+b^2}$}

\begin{corrige}
	On a $\big\|\vv{v}(\ptM)\big\|=\sqrt{\dot{x}^2+\dot{y}^2+\dot{z}^2}=\sqrt{\left(a\omega\right)^2\big(\sin(\omega t)^2+\cos(\omega t)^2\big)+b^2}=\sqrt{\left(a\omega\right)^2+b^2}$.
\end{corrige}

% ************************************** %


% ^^^^^^^^^^^^^^^^^^^^^^^^^^^^^^^^^^^^^^ %
%               Question                 %
% ^^^^^^^^^^^^^^^^^^^^^^^^^^^^^^^^^^^^^^ %

\begin{enonce}
	Déterminer l'accélération $\vv{a}(\ptM)$ dans la base cartésienne
\end{enonce}

\reponse{$-a\omega^2\left(\cos(\omega t) \vv{e_x}+\sin(\omega t) \vv{e_y}\right)$}

\begin{corrige}
	On a $\vv{a}(\ptM)=\ddot{x}\, \vv{e_x}+\ddot{y}\, \vv{e_y}+\ddot{z}\, \vv{e_z}=
	-a\omega^2\big(\cos(\omega t) \vv{e_x}+\sin(\omega t) \vv{e_y}\big)$.
\end{corrige}

% ************************************** %


% ^^^^^^^^^^^^^^^^^^^^^^^^^^^^^^^^^^^^^^ %
%               Question                 %
% ^^^^^^^^^^^^^^^^^^^^^^^^^^^^^^^^^^^^^^ %

\begin{enonce}
	Déterminer la norme de l'accélération  
\end{enonce}

\reponse{$a\omega^2$}

\begin{corrige}
	On a $\big\|\vv{a}(\ptM)\big\|=\sqrt{\ddot{x}^2+\ddot{y}^2+\ddot{z}^2}=a\omega^2$.
\end{corrige}

% ************************************** %


% =============================================================================== %
\finEntrainement
% =============================================================================== %





% ********************************************************************* %
%                           ENTRAÎNEMENT                                % 
% ********************************************************************* %
% ================ Métadonnées sur l'entraînement ===================== %

\titreEntrainementFacultatif	{Dérivation des vecteurs unitaires de la base polaire}
\hauteurLargeurCadreReponse		{8mm}{5cm}
\dureeResolutionFacultative		{1} % 1, 2, 3 ou 4
\typeEntrainement				{Newton} % Newton ou QCM ou AN ou vide (exercice "normal")
\nombreColonnesQuestions		{1} % vide, 1, 2, 3, etc.
\avecPlusieursQuestions			{Y} % Y ou N
\initialisationEntrainement
% ===================================================================== %



% mmmmmmmmmmmmmmmmmmmmmmmmmmmmmmmmmmmmmmmmmmmmmmmmmmmmmmmmm %
% 		  Insertion d'une image en regard d'un texte
% mmmmmmmmmmmmmmmmmmmmmmmmmmmmmmmmmmmmmmmmmmmmmmmmmmmmmmmmm %
% --------------------------------------------------------- %
\pourcentageDeLaPartieAGauche   {0.6}
% --------------------------------------------------------- %
%                                                           %
% -------------------- Partie à gauche -------------------- %
%                                                           %
\initialisationPartieGauche % 
%                                                           %
%                                                           %
On considère un point $\ptM(t)$ en mouvement dans le plan $(x \ptO y)$. 

On note $r(t)$ et $\theta(t)$ les coordonnées de $\ptM(t)$ 
dans le repère polaire $(\ptO, \vv{e_r}, \vv{e_\theta})$.
% --------------------------------------------------------- %
%                                                           %
% -------------------- Partie à droite -------------------- %
%                                                           %
\initialisationPartieDroite %
%                                                           %
%                                                           %
\begin{center}
	\subimport{_images/}{fiche_sup_1.9_ECA.tex}
\end{center}
% --------------------------------------------------------- %
\finalisationDuPartageDePage %					
% --------------------------------------------------------- %


% =============================================================================== %
\debutEntrainement
% =============================================================================== %


% ^^^^^^^^^^^^^^^^^^^^^^^^^^^^^^^^^^^^^^ %
%               Question                 %
% ^^^^^^^^^^^^^^^^^^^^^^^^^^^^^^^^^^^^^^ %
\begin{enonce}
	Exprimer le vecteur  $\vv{e_r}$ dans la base 
	cartésienne $(\ptO, \vv{e_x}, \vv{e_y})$
\end{enonce}

\reponse{$\cos \theta \vv{e_x} + \sin \theta \vv{e_y}$}

% ************************************** %


% ^^^^^^^^^^^^^^^^^^^^^^^^^^^^^^^^^^^^^^ %
%               Question                 %
% ^^^^^^^^^^^^^^^^^^^^^^^^^^^^^^^^^^^^^^ %

\begin{enonce}
En déduire la dérivée  $\diff{\vv{e_r}}{t}$ dans la base 
cartésienne $(\ptO, \vv{e_x}, \vv{e_y})$
\end{enonce}

\reponse{$\diff{\vv{e_r}}{t} = \dot{\theta}\left(-\sin\theta \vv{e_x} + 
\cos\theta\vv{e_y}\right)$}

% ************************************** %


% ^^^^^^^^^^^^^^^^^^^^^^^^^^^^^^^^^^^^^^ %
%               Question                 %
% ^^^^^^^^^^^^^^^^^^^^^^^^^^^^^^^^^^^^^^ %

\begin{enonce}
	Exprimer le vecteur   $\vv{e_x}$ dans la base 
	polaire $(\ptO, \vv{e_r}, \vv{e_\theta})$
\end{enonce}

\reponse{$\vv{e_x}=\cos \theta \vv{e_r} -\sin\theta \vv{e_\theta}$}

% ************************************** %


% ^^^^^^^^^^^^^^^^^^^^^^^^^^^^^^^^^^^^^^ %
%               Question                 %
% ^^^^^^^^^^^^^^^^^^^^^^^^^^^^^^^^^^^^^^ %

\begin{enonce}
	Exprimer le vecteur  $\vv{e_y}$ dans la base 
	polaire $(\ptO, \vv{e_r}, \vv{e_\theta})$
\end{enonce}

\reponse{$\vv{e_y}=\sin \theta \vv{e_r} +\cos\theta \vv{e_\theta}$}

% ************************************** %



% ^^^^^^^^^^^^^^^^^^^^^^^^^^^^^^^^^^^^^^ %
%               Question                 %
% ^^^^^^^^^^^^^^^^^^^^^^^^^^^^^^^^^^^^^^ %

\begin{enonce}
	En déduire l'expression de la dérivée  $\diff{\vv{e_r}}{t}$ dans la base 
	polaire
\end{enonce}

\reponse{$\diff{\vv{e_r}}{t} = \dot{\theta}\vv{e_\theta}$}

% ************************************** %

% =============================================================================== %
\finEntrainement
% =============================================================================== %






\newpage




% ********************************************************************* %
%                           ENTRAÎNEMENT                                % 
% ********************************************************************* %
% ================ Métadonnées sur l'entraînement ===================== %

\titreEntrainementFacultatif	{Calcul d'une vitesse en coordonnées polaires}
\hauteurLargeurCadreReponse		{8mm}{4cm}
\dureeResolutionFacultative		{1} % 1, 2, 3 ou 4
\typeEntrainement				{} % Newton ou QCM ou AN ou vide (exercice "normal")
\basiqueEtTransversal			{Y}
\nombreColonnesQuestions		{1} % vide, 1, 2, 3, etc.
\avecPlusieursQuestions			{Y} % Y ou N
\initialisationEntrainement
% ===================================================================== %


On considère un point $\ptM$ dont les coordonnées polaires sont 
$\begin{cases}
	r(t) =a\times t\\
	\theta(t) = b\times t^2.
\end{cases}$
\newline
La vitesse en coordonnées polaires s'écrit : 
\[
\vv{v}(\ptM)=\dot{r} \, \vv{e_r}+ r\dot{\theta}\, \vv{e_\theta}, 
\]
où $\dot{r}\, \vv{e_r}$ est appelée \emph{vitesse radiale} et $r\dot{\theta}\, \vv{e_\theta}$  \emph{vitesse orthoradiale}. 

% =============================================================================== %
\debutEntrainement
% =============================================================================== %


% ^^^^^^^^^^^^^^^^^^^^^^^^^^^^^^^^^^^^^^ %
%               Question                 %
% ^^^^^^^^^^^^^^^^^^^^^^^^^^^^^^^^^^^^^^ %

\begin{enonce}
	Déterminer la dimension de $a$
\end{enonce}

\reponse{$\frac{L}{T}$}

\begin{corrige}
	On a $a=\frac{r}{t}$. Ainsi, $a$ est homogène à une longueur sur un temps.
\end{corrige}

% ************************************** %


% ^^^^^^^^^^^^^^^^^^^^^^^^^^^^^^^^^^^^^^ %
%               Question                 %
% ^^^^^^^^^^^^^^^^^^^^^^^^^^^^^^^^^^^^^^ %

\begin{enonce}
	Déterminer la dimension de $b$
\end{enonce}

\reponse{$\frac{1}{T^2}$}

\begin{corrige}
	On a $b=\frac{\theta}{t^2}$. Ainsi, $b$ est homogène à un angle sur un temps au carré. Comme un angle est une grandeur sans dimension, on a bien le résultat donné. 
\end{corrige}

% ************************************** %


% ^^^^^^^^^^^^^^^^^^^^^^^^^^^^^^^^^^^^^^ %
%               Question                 %
% ^^^^^^^^^^^^^^^^^^^^^^^^^^^^^^^^^^^^^^ %

\begin{enonce}
	Déterminer la vitesse radiale en fonction de $a$
\end{enonce}

\reponse{$a\vv{e_r}$}

\begin{corrige}
	La vitesse radiale est $\vv{v}(\ptM)_r=\dot{r} \vv{e_r}=a \vv{e_r}$.
\end{corrige}

% ************************************** %


% ^^^^^^^^^^^^^^^^^^^^^^^^^^^^^^^^^^^^^^ %
%               Question                 %
% ^^^^^^^^^^^^^^^^^^^^^^^^^^^^^^^^^^^^^^ %

\begin{enonce}
	Déterminer la vitesse orthoradiale  en fonction de $a$, $b$ et $t$
\end{enonce}

\reponse{$2abt^2 \vv{e_\theta}$}

\begin{corrige}
	La vitesse orthoradiale est $\vv{v}(\ptM)_\theta=r\dot{\theta} \vv{e_\theta}=2abt^2 \vv{e_\theta}$
\end{corrige}

% ************************************** %


% ^^^^^^^^^^^^^^^^^^^^^^^^^^^^^^^^^^^^^^ %
%               Question                 %
% ^^^^^^^^^^^^^^^^^^^^^^^^^^^^^^^^^^^^^^ %

\begin{enonce}
	En déduire l'expression de $\vv{v}(\ptM)$ 
\end{enonce}

\reponse{$a\vv{e_r}+2abt^2 \vv{e_\theta}$}

\begin{corrige}
	On a $\vv{v}(\ptM)=a\vv{e_r}+2abt^2 \vv{e_\theta}$.
\end{corrige}

% ************************************** %

% =============================================================================== %
\finEntrainement
% =============================================================================== %






% ********************************************************************* %
%                           ENTRAÎNEMENT                                % 
% ********************************************************************* %
% ================ Métadonnées sur l'entraînement ===================== %

\titreEntrainementFacultatif	{Mouvement en spirale}
\hauteurLargeurCadreReponse		{8mm}{4cm}
\dureeResolutionFacultative		{2} % 1, 2, 3 ou 4
\typeEntrainement				{} % Newton ou QCM ou AN ou vide (exercice "normal")
\nombreColonnesQuestions		{1} % vide, 1, 2, 3, etc.
\avecPlusieursQuestions			{Y} % Y ou N
\initialisationEntrainement
% ===================================================================== %


Un point $\ptM(t)$ décrit une trajectoire en forme de spirale. 
Dans le repère polaire $(\ptO, \vv{e_r}, \vv{e_\theta})$, les coordonnées de $\ptM(t)$ sont :
\[
\begin{cases}
	r(t) = r_0\, \eEuler^{-t/\tau}\\
	\theta(t)  = \omega t 
\end{cases}
\]
où $r_0$, $\tau$ et $\omega$ sont des constantes positives.


% =============================================================================== %
\debutEntrainement
% =============================================================================== %





% ^^^^^^^^^^^^^^^^^^^^^^^^^^^^^^^^^^^^^^ %
%               Question                 %
% ^^^^^^^^^^^^^^^^^^^^^^^^^^^^^^^^^^^^^^ %

\begin{enonce}
	Déterminer la vitesse $\vv{v}(\ptM)$ en coordonnées polaires. \minusMedskip
	
	{\itshape On pourra utiliser la formule donnée dans l'entraînement précédent}
\end{enonce}

\reponse{$r_0\eEuler^{-t/\tau}\left(-\frac{1}{\tau}\vv{e_r}+\omega\vv{e_\theta} \right)$}

\begin{corrige}
	On a $ \vv{v}(\ptM)=\dot{r} \vv{e_r}+ r\dot{\theta}\vv{e_\theta}=r_0 \eEuler^{-t/\tau}\left(-\frac{1}{\tau}\vv{e_r}+\omega\vv{e_\theta} \right)$.
\end{corrige}	

% ************************************** %


% =============================================================================== %
\pauseEntrainement
% =============================================================================== %

L'accélération en coordonnées polaires s'écrit : 
\[
\vv{a}(\ptM)=\left(\ddot{r}-r\dot{\theta}^2\right) \vv{e_r}+\left(2\dot{r}\dot{\theta}+r\ddot{\theta}\right) \vv{e_\theta}.
\]
\vspace{-8mm}

% =============================================================================== %
\repriseEntrainement
% =============================================================================== %



% ^^^^^^^^^^^^^^^^^^^^^^^^^^^^^^^^^^^^^^ %
%               Question                 %
% ^^^^^^^^^^^^^^^^^^^^^^^^^^^^^^^^^^^^^^ %

\begin{enonce}
	Déterminer l'accélération $\vv{a}(\ptM)$
\end{enonce}

\reponse{$r_0 \eEuler^{-t/\tau}  \left(\left(\frac{1}{\tau^2}-\omega^2\right)\vv{e_r}-\left(2\frac{\omega}{\tau}\right)\vv{e_\theta}\right)$}

\begin{corrige}
	On a $\vv{a}(\ptM)=r_0 \eEuler^{-t/\tau}  \left(\left(\frac{1}{\tau^2}-\omega^2\right)\vv{e_r}-\left(2\frac{\omega}{\tau}\right)\vv{e_\theta}\right)$.
\end{corrige}	

% ************************************** %


% =============================================================================== %
\pauseEntrainement
% =============================================================================== %

On donne les valeurs suivantes : 
$\omega=\SI{4,78}{tour.min^{-1}}$,  $\tau=\SI{2,0}{s}$ et $r_0=\SI{4,0}{cm}$. 

% =============================================================================== %
\repriseEntrainement
% =============================================================================== %


% ^^^^^^^^^^^^^^^^^^^^^^^^^^^^^^^^^^^^^^ %
%               Question                 %
% ^^^^^^^^^^^^^^^^^^^^^^^^^^^^^^^^^^^^^^ %

\begin{enonce}
	Dans ces conditions, l'accélération est-elle radiale ou orthoradiale ?
\end{enonce}

\reponse{orthoradiale}

\begin{corrige}
	On a $\omega=\SI{4,78}{tour.min^{-1}}=\frac{\np{4,78}\times 2\pi\ \si{rad}}{\SI{60}{s}}=\SI{0,5}{rad.s^{-1}}$ et 
	$\frac{1}{\tau^2}-\omega^2=\left(\frac{1}{\np{0,5}^2}-\omega^2 \right)=\SI{0}{s^{-2}}$. 
	
	Ainsi, on a 
	$\vv{a}(\ptM,t)=-2r_0 \eEuler^{-t/\tau} \vv{e_\theta}$.	L'accélération est donc orthoradiale. 
\end{corrige}	

% ************************************** %


% ^^^^^^^^^^^^^^^^^^^^^^^^^^^^^^^^^^^^^^ %
%               Question                 %
% ^^^^^^^^^^^^^^^^^^^^^^^^^^^^^^^^^^^^^^ %

\begin{enonce}
	Le mouvement de $\ptM$ est-il accéléré ou décéléré ?
\end{enonce}

\reponse{décéléré}

\begin{corrige}
	On a $\vv{a}(\ptM,t)\cdot \vv{e_\theta} =-2r_0 \eEuler^{-t/\tau} <0$.
	Le mouvement est donc décéléré. 
\end{corrige}	

% ************************************** %


% ^^^^^^^^^^^^^^^^^^^^^^^^^^^^^^^^^^^^^^ %
%               Question                 %
% ^^^^^^^^^^^^^^^^^^^^^^^^^^^^^^^^^^^^^^ %

\hauteurLargeurCadreReponse		{8mm}{6cm}

\begin{enonce}
	Déterminer l'équation polaire de la trajectoire de $\ptM$
\end{enonce}

\reponse{$r = r_0\eEuler^{-\theta}$}

\begin{corrige}
	On a $r = r_0e^{-t/\tau}$ et $t=\frac{\theta}{\omega}$. Donc, on a 
	$r = r_0\eEuler^{-\theta/(\omega\times \tau)} = r_0\eEuler^{-\theta}$ car $\omega \tau = 1$.
\end{corrige}

% ************************************** %


% =============================================================================== %
\finEntrainement
% =============================================================================== %



\newpage



% •••••••••••••••••••••••••••••••••••••••••••••••••••••••••••••••••••••••••••• %
% •••••••••••••••••••••••••••••••••••••••••••••••••••••••••••••••••••••••••••• %
\sectionFicheEntrainement{Étude de quelques mouvements}
% •••••••••••••••••••••••••••••••••••••••••••••••••••••••••••••••••••••••••••• %
% •••••••••••••••••••••••••••••••••••••••••••••••••••••••••••••••••••••••••••• %




% ********************************************************************* %
%                           ENTRAÎNEMENT                                % 
% ********************************************************************* %
% ================ Métadonnées sur l'entraînement ===================== %

\titreEntrainementFacultatif	{Collision sur plan incliné}
\hauteurLargeurCadreReponse		{8mm}{4cm}
\dureeResolutionFacultative		{4} % 1, 2, 3 ou 4
\typeEntrainement				{Newton} % Newton ou QCM ou AN ou vide (exercice "normal")
\nombreColonnesQuestions		{1} % vide, 1, 2, 3, etc.
\avecPlusieursQuestions			{Y} % Y ou N
\initialisationEntrainement
% ===================================================================== %



% mmmmmmmmmmmmmmmmmmmmmmmmmmmmmmmmmmmmmmmmmmmmmmmmmmmmmmmmm %
% 		  Insertion d'une image en regard d'un texte
% mmmmmmmmmmmmmmmmmmmmmmmmmmmmmmmmmmmmmmmmmmmmmmmmmmmmmmmmm %
% --------------------------------------------------------- %
\pourcentageDeLaPartieAGauche   {0.65}
% --------------------------------------------------------- %
%                                                           %
% -------------------- Partie à gauche -------------------- %
%                                                           %
\initialisationPartieGauche % 
%                                                           %
%                                                           %
Deux billes évoluent sur un plan incliné faisant un angle $\alpha=\SI{20}{\degree}$ par rapport à l'horizontale. 

À $t=0$, elles sont distantes d'une longueur $L$.

\begin{itemize}
\item
La bille $\ptA$ possède une vitesse initiale 
$v_0 \vv{e_{x'}}$.

Son accélération $\vv{a}(A) =-a\vv{e_{x'}}$ est constante au cours du temps. 

Nous noterons  
$v_A(t)  \vv{e_{x'}}$ sa vitesse à l'instant $t$.  
\smallskip

\item
La bille $\ptB$ quant à elle, n'a pas de vitesse initiale mais possède une accélération constante $\vv{a}(B) =a\vv{e_{x'}}$. 

Nous noterons  $v_B(t)  \vv{e_{x'}}$ sa vitesse à l'instant $t$. 
\end{itemize}


% --------------------------------------------------------- %
%                                                           %
% -------------------- Partie à droite -------------------- %
%                                                           %
\initialisationPartieDroite %
%                                                           %
%                                                           %
\begin{center}
	\subimport{_images/}{fiche_sup_choc_CBD_v2.tex}
\end{center}
% --------------------------------------------------------- %
\finalisationDuPartageDePage %					
% --------------------------------------------------------- %

\medskip
On donne  $a=\SI{3,35}{m.s^{-2}}$ et $v_0=\SI{3}{m.s^{-1}}$.


% =============================================================================== %
\debutEntrainement
% =============================================================================== %


% ^^^^^^^^^^^^^^^^^^^^^^^^^^^^^^^^^^^^^^ %
%               Question                 %
% ^^^^^^^^^^^^^^^^^^^^^^^^^^^^^^^^^^^^^^ %

\begin{enonce}
Exprimer $v_A(t)$  en fonction $a$, $t$ et $v_0$ 
\end{enonce}

\reponse{$-at+v_0$}

\begin{corrige}
	On a $\vv{a}(A) = 
	\diff{\vv{v}(A)}{t}$.  
	En projetant sur l'axe $(0,\vv{e_{x'}})$, on obtient : $-a = 
	\diff{v_A}{t} $. Puis, en calculant
	\[
	\int^{v_A(t)}_{v_0}	\d {v_A}=\int^t_0 -a  \d t,
	\]
	on obtient $v_A(t)=-at+v_0$.
\end{corrige}

% ************************************** %


% ^^^^^^^^^^^^^^^^^^^^^^^^^^^^^^^^^^^^^^ %
%               Question                 %
% ^^^^^^^^^^^^^^^^^^^^^^^^^^^^^^^^^^^^^^ %

\begin{enonce}
	Exprimer $v_B(t)$  en fonction $a$ et $t$
\end{enonce}

\reponse{$at$}

\begin{corrige}
	On a $\vv{a}(B) = 
	\diff{\vv{v}(B)}{t}$. 
	En projetant sur l'axe $(0,\vv{e_{x'}})$, on obtient : $a = 
	\diff{v_B}{t}$. Puis, en calculant
	\[
	\int^{v_B(t)}_{0} \d {v_B}=\int^t_0 a  \d t,
	\]
	on obtient  $v_B(t)=at$.
\end{corrige}

% ************************************** %


% ^^^^^^^^^^^^^^^^^^^^^^^^^^^^^^^^^^^^^^ %
%               Question                 %
% ^^^^^^^^^^^^^^^^^^^^^^^^^^^^^^^^^^^^^^ %

\begin{enonce}
	Déterminer la position $x'_A$ de $\ptA$ en fonction du temps
\end{enonce}

\reponse{$-\frac{1}{2}at^2+v_0t$}

\begin{corrige}
	Sur l'axe $(0,\vv{e_{x'}})$,  on a : $v_A(t) = 
	\diff{x'_A}{t} $. 
	Donc, on a
	$\int^{x'_A(t)}_{0}	\d {x'_A}=\int^t_0 v_A  \d t=\int^t_0 \left(-at+v_0\right)  \d t$. 
	
	Donc, 
	on a  $x'_A(t)=-\frac{1}{2}at^2+v_0t$.
\end{corrige}

% ************************************** %


% ^^^^^^^^^^^^^^^^^^^^^^^^^^^^^^^^^^^^^^ %
%               Question                 %
% ^^^^^^^^^^^^^^^^^^^^^^^^^^^^^^^^^^^^^^ %

\begin{enonce}
	Déterminer la position $x'_B$ de $\ptB$ en fonction du temps	
\end{enonce}

\reponse{$\frac{1}{2}at^2+L$}

\begin{corrige}
	Sur l'axe $(0,\vv{e_{x'}})$ nous avons : $v_B(t) = 
	\diff{x'_B}{t} $. Donc, on a
	$\int^{x_B(t)}_{L}	\d {x'_B}=\int^t_0 v_B  \d t=\int^t_0 at  \d t$.
	
		Donc, 
	on a  $x'_B(t)=\frac{1}{2}at^2+L$.
\end{corrige}

% ************************************** %


% ^^^^^^^^^^^^^^^^^^^^^^^^^^^^^^^^^^^^^^ %
%               Question                 %
% ^^^^^^^^^^^^^^^^^^^^^^^^^^^^^^^^^^^^^^ %

\begin{enonce}
	Déterminer la distance $L$ minimale (en $\si{cm}$) pour qu'une collision puisse avoir lieu.
	
\end{enonce}

\reponse{$\SI{67}{cm}$}

\begin{corrige}
	Nous observerons une collision à la date $t_1$ si $x'_A(t_1)=x'_B(t_1)$ donc si
	$-\frac{1}{2}at_1^2+v_0t_1=\frac{1}{2}at_1^2+L$.
	
	Donc, $t_1$ doit être une solution réelle positive de l'équation : 
	\[
	t_1^2-\frac{v_0}{a}t_1+\frac{L}{a}=0,
	\]
	ce qui impose une valeur positive pour son discriminant $\Delta= \left(\frac{v_0}{a}\right)^2-4\frac{L}{a} \geq 0$. Donc, on doit avoir $L \leq \frac{v_0^2}{4a}$.
	
	Après application numérique, on trouve que la distance $L$ doit vérifier $L\leq \SI{67}{cm}$.
\end{corrige}

% ************************************** %

% =============================================================================== %
\finEntrainement
% =============================================================================== %






% ********************************************************************* %
%                           ENTRAÎNEMENT                                % 
% ********************************************************************* %
% ================ Métadonnées sur l'entraînement ===================== %

\titreEntrainementFacultatif	{Chute libre}
\hauteurLargeurCadreReponse		{8mm}{4cm}
\dureeResolutionFacultative		{2} % 1, 2, 3 ou 4
\typeEntrainement				{Newton} % Newton ou QCM ou AN ou vide (exercice "normal")
\nombreColonnesQuestions		{1} % vide, 1, 2, 3, etc.
\avecPlusieursQuestions			{Y} % Y ou N
\initialisationEntrainement
% ===================================================================== %



On considère le point $\ptM$ de masse $m$ et de coordonnées ($x,y,z$) dans 
la base cartésienne $(\ptO,\vv{e_x},\vv{e_y},\vv{e_z})$. 

Il est lancé avec la 
vitesse $\vv{v_0} = v_{0x}\vv{e_x} + v_{0z}\vv{e_z}$ à partir de l'origine $\ptO$ du repère dans le champ de pesanteur uniforme $\vv{g} = -g\vv{e_z}$.

Tout frottement étant négligé, l'accélération de $\ptM$ est égale à $\vv{g}$ à tout instant. 


% =============================================================================== %
\debutEntrainement
% =============================================================================== %


% ^^^^^^^^^^^^^^^^^^^^^^^^^^^^^^^^^^^^^^ %
%               Question                 %
% ^^^^^^^^^^^^^^^^^^^^^^^^^^^^^^^^^^^^^^ %

\begin{enonce}
	Exprimer $x(t)$ en fonction de $v_{0x}$ et $t$
\end{enonce}

\reponse{$v_{0x}t$}
	
\begin{corrigeNewpage}
 On a $\vv{a} = 
	\diff{\vv{v}}{t} = \vv{g}$.
	En projetant, nous obtenons 
	\[
	\begin{cases}
		\diff{v_x}{t} = 0\\[2mm]
		\diff{v_z}{t} = -g.
	\end{cases}
	\]
Donc, on a $v_x = \mathrm{C^{te}}=v_{0x}$. En intégrant une deuxième fois,  vu que $\ptM$ 
est initialement en $\ptO$, on obtient : $x(t) = v_{0x}t$.
\end{corrigeNewpage}

% ************************************** %


% ^^^^^^^^^^^^^^^^^^^^^^^^^^^^^^^^^^^^^^ %
%               Question                 %
% ^^^^^^^^^^^^^^^^^^^^^^^^^^^^^^^^^^^^^^ %

\begin{enonce}
	Exprimer $z(t)$ en fonction de $v_{0z}$, $g$ et $t$
\end{enonce}

\reponse{$-\dfrac{1}{2}gt^2+v_{0z}t$}
	
\begin{corrige}
 On a $\vv{a} = 
	\diff{\vv{v}}{t} = \vv{g}$.
	En projetant, nous obtenons 
	\[
	\begin{cases}
		\diff{v_x}{t} = 0\\[2mm]
		\diff{v_z}{t} = -g.
	\end{cases}
	\]
Donc, en intégrant,  on a $\int^{v_z(t)}_{v_{0z}}\d{v_z}=\int^t_0 -g \cdot \d t$ donc $v_z=-gt+v_{0z}$.
En intégrant une deuxième fois,  vu que $\ptM$ 
est initialement en $\ptO$, on obtient : 
\[
	z(t) = -\dfrac{1}{2}gt^2 + v_{0z}t.
\]
\end{corrige}

% ************************************** %


% ^^^^^^^^^^^^^^^^^^^^^^^^^^^^^^^^^^^^^^ %
%               Question                 %
% ^^^^^^^^^^^^^^^^^^^^^^^^^^^^^^^^^^^^^^ %

\smallskip

\begin{enonce}
	En déduire l'équation cartésienne de la trajectoire $z$ en fonction de $x$, \minusBigskip
	\par c'est-à-dire une relation entre $x(t)$ et $z(t)$
\end{enonce}

\reponse{$z=-\dfrac{g}{2v_{0x}^2}x^2+\dfrac{v_{0z}}{v_{0x}}x$}

\begin{corrige}
À partir de l'expression de $x(t)$ on peut écrire $t=x/v_{0x}$. On remplace $t$ par cette expression dans $z$:
\[
z =-\dfrac{1}{2}g(x/v_{0x})^2 + v_{0z} x/v_{0x}.
\]
Finalement, on trouve l'équation $z=-\dfrac{g}{2v_{0x}^2}x^2+\dfrac{v_{0z}}{v_{0x}}x$.
\end{corrige}

% ************************************** %

% =============================================================================== %
\finEntrainement
% =============================================================================== %



\newpage



% ********************************************************************* %
%                           ENTRAÎNEMENT                                % 
% ********************************************************************* %
% ================ Métadonnées sur l'entraînement ===================== %

\titreEntrainementFacultatif	{Pauvre gazelle}
\hauteurLargeurCadreReponse		{8mm}{3cm}
\dureeResolutionFacultative		{2} % 1, 2, 3 ou 4
\typeEntrainement				{Newton} % Newton ou QCM ou AN ou vide (exercice "normal")
\nombreColonnesQuestions		{1} % vide, 1, 2, 3, etc.
\avecPlusieursQuestions			{Y} % Y ou N
\initialisationEntrainement
% ===================================================================== %

Un lion chasse une gazelle. Il court à la vitesse constante de 
$\SI{5,0}{m.s^{-1}}$. La gazelle aperçoit  le lion quand il est à \SI{10}{m} de 
distance. Elle se met alors en fuite en accélérant à $\SI{2,0}{m.s^{-2}}$. 
Pour rattraper la gazelle, le lion se met aussi à accélérer au même 
instant à $\SI{3,0}{m.s^{-2}}$.


% =============================================================================== %
\debutEntrainement
% =============================================================================== %

% ^^^^^^^^^^^^^^^^^^^^^^^^^^^^^^^^^^^^^^ %
%               Question                 %
% ^^^^^^^^^^^^^^^^^^^^^^^^^^^^^^^^^^^^^^ %

\begin{enonce}
	Combien de temps mettra le lion à rattraper la gazelle ?
\end{enonce}

\reponse{$\SI{1,7}{s}$}

\begin{corrige}
	On suppose que le lion et la gazelle se déplacent en ligne droite sur 
	l'axe $(\ptO x)$. On prend l'origine des temps au moment où la gazelle 
	aperçoit le lion et l'origine de l'axe $(\ptO x)$ à la position du lion quand 
	la gazelle l'aperçoit. 
	
	On intègre deux fois pour avoir la position du 
	lion $x_L$ puis celle de la gazelle $x_G$ en fonction de temps: 
	\[
	\begin{cases}
		x_L(t) = v_0t + \dfrac{1}{2}a_Lt^2\\[2mm]
		x_G(t) = d_0 + \dfrac{1}{2}a_Gt^2.
	\end{cases}
	\]
avec $v_0 =\SI{5,0}{m.s^{-1}}$, $a_L=\SI{3,0}{m.s^{-2}}$, $a_G=\SI{2,0}{m.s^{-2}}$ et $d_0=\SI{10}{m}$. 

	Puis, on égalise ces 
	deux positions pour déterminer le temps $t_1$ où le lion attrape la 
	gazelle.  On obtient une équation du second degré sur $t_1$: 
	\[
	\dfrac{a_L-a_G}{2}t_1^2 + v_0 t_1 - d_0 = 0. \tag{$\ast$}
	\]
	On résout cette équation du second degré qui admet deux racines réelles 
	dont l'une est négative. Le temps cherché est la racine 
	positive : c'est $t_1 = \dfrac{-v_0 + \sqrt{\Delta}}{a_L-a_G}$ où $\Delta = 
	v_0^2 + 2d_0(a_L-a_G)$ est le discriminant de l'équation ($\ast$).
	
	On trouve finalement $t_1 = \SI{1,7}{s}$.
\end{corrige}

% ************************************** %


% ^^^^^^^^^^^^^^^^^^^^^^^^^^^^^^^^^^^^^^ %
%               Question                 %
% ^^^^^^^^^^^^^^^^^^^^^^^^^^^^^^^^^^^^^^ %

\begin{enonce}
	Quelle distance aura parcourue la gazelle avant de se faire dévorer ?
\end{enonce}

\reponse{$\SI{2,9}{m}$}

\begin{corrige}
	La gazelle aura parcouru la distance $d = \dfrac{1}{2}a_G t_1^2$ avec 
	$a_G = \SI{2,0}{m.s^{-2}} $ et $t_1 = \SI{1,7}{s}$ le temps mis par le lion 
	pour rattraper la gazelle. Finalement, on trouve  $d =\SI{2,9}{m}$.
\end{corrige}

% ************************************** %

% =============================================================================== %
\finEntrainement
% =============================================================================== %







% ---------------------------------------------------------------------------- %
%                       Affichage des réponses mélangées                       %
% ---------------------------------------------------------------------------- %
\afficheReponsesMelangees
% ---------------------------------------------------------------------------- %

%%%%%%%%%%%%%%%%%%%%%%%%%%%%%%%%%%%%%%%%%%%%%%%%%%%%%%%%%%%%%%%%%%%%%%%%%%%%%%%%%%%%%%%%%%%%%%
\finFicheEntrainement                                                                        %
%%%%%%%%%%%%%%%%%%%%%%%%%%%%%%%%%%%%%%%%%%%%%%%%%%%%%%%%%%%%%%%%%%%%%%%%%%%%%%%%%%%%%%%%%%%%%%


% ======================================================= % 
% ============== Gestion de la compilation ============== %
% ======================================================= % 
\ifdefined\mainIsLoaded\else
\printReponsesEtCorriges
\end{document}\fi
% ======================================================= % 
% ======================================================= % 

